\section{Introduction}

What remains random in a shuffle after the paths of many cards have
been observed? For the Thorp shuffle, the answer is especially concrete.
On a pair containing two unobserved positions, the conditional masses
of an additional card are averaged. On a pair containing one observed
position, the remaining mass is transported to the uniquely available
output. Observing more paths therefore changes an averaging process
into a random mixture of averaging and transport. This paper estimates
the information retained by that conditional process.

The positions are the vertices of $V=\mathbb F_2^d$, with $n=2^d$.
One physical Thorp shuffle independently switches each pair in one
coordinate direction and rotates the coordinates. Undoing these
rotations gives a process $(\Pi_t)$ that uses the coordinate directions
in a fixed cyclic order. A sweep consists of $d$ such layers, hence
$d$ physical shuffles. We label cards by their initial positions;
$\Pi_t(a)$ is the position of card $a$ at time $t$. All starting
decks considered below are deterministic. Total variation has the
normalization $\|\mu-\nu\|_{\TV}=\frac12\sum_x|\mu(x)-\nu(x)|$.
Write $q_d$ for the law of one physical shuffle on $S_n$ and $U_{S_n}$
for the uniform law. The threshold-$1/4$ mixing time is
\[
 t_{\mathrm{mix}}(d)=\min\{t\in\mathbb Z_{\ge0}:
          \|q_d^{*t}-U_{S_n}\|_{\TV}\le1/4\}.
\]
Convolution represents independent successive shuffles. Relabeling
the deterministic starting deck leaves its distance from uniform unchanged.

Thorp introduced the shuffle in a study of nonuniform human shuffling
and its consequences for Faro~\cite{thorp,morris44}. Its simple description
conceals a difficult full-deck mixing problem: the transitions are
nonreversible, and uniformity of individual card positions does not
control their joint order. Morris's first polynomial bound for $n=2^d$,
circulated in 2005 and published in 2008, gave $O(d^{44})$ physical
shuffles~\cite{morris44}. His proof combined evolving sets with a
chameleon process for a forward--backward version of the shuffle.
Montenegro and Tetali retained Morris's contraction estimate and
sharpened the mixing analysis using spectral profiles and evolving
sets, obtaining $O(d^{29})$ in 2006
\cite[Theorems~6.14--6.15]{montenegro-tetali2006}.

Morris subsequently developed an entropy method that reduces mixing
estimates to local interactions between pairs of cards. It yielded
$O((\log n)^4)$ physical shuffles for every even $n$~\cite{morris4},
followed by $O(d^3)$ for $n=2^d$~\cite{morris3}. These arguments
decompose the entropy of a permutation into conditional one-card
contributions and analyze the reverse shuffle. They make explicit
why information remaining after other cards are exposed is central
to the full-deck problem.

Ordered Thorp marginals also arise in the small-domain enciphering
analysis of Morris, Rogaway and Stegers~\cite[Theorem~1]{mrs}.
For each fixed $0<\varepsilon<1$, their estimate gives
$O_\varepsilon(d)$ physical shuffles for specified lists of at most
$n^{1-\varepsilon}$ cards.
Czumaj and V\"ocking~\cite{CzumajVocking2014} reached lists containing
any fixed fraction less than one of the deck: a non-Markovian coupling
on butterfly networks gives $O((\log n)^2)$ physical shuffles.
Their full-permutation algorithm adds empty cards and removes them
afterwards; its chain therefore differs from the unchanged labeled
deck considered here.

A closely related conditional-energy strategy appears in the
swap-or-not analysis of Hoang, Morris and Rogaway
\cite[Section~3, Theorem~3 and Lemma~4]{hmr2014}.
They control a centered squared error for the next card and sum the
resulting conditional errors along an ordered list. Their matching
is selected by a fresh uniform group key at each round, and their
contraction is averaged over those keys. Here the coordinate directions
are prescribed cyclically, while the switch coins remain random.
We condition explicitly on preceding cards' complete paths. The
resulting transport-and-average flow admits an exact one-sweep memory
identity, and independent coordinate blocks control its remaining
noise. Thus the common sequential comparison is retained, while the
contraction comes from the geometry of the deterministic schedule.

\paragraph{Conditional information and the first mixing result.}
For a fixed set $B$ of observed labels, let $\mathcal H_t$ record its
paths through time $t$, and let $V_t=V\setminus\Pi_t(B)$ be the
available positions. For a further label $a\notin B$, the basic object is
\[
 f_t(x)=\Pr\{\Pi_t(a)=x\mid\mathcal H_t\}
             -\frac{\one_{V_t}(x)}{|V_t|}.
\]
It is a sum-zero vector supported on $V_t$. Its energy
$\E\|f_t\|_2^2$ controls the expected error in the conditional law
of $a$. Revealing additional cards successively converts this scalar
estimate into information about an ordered list. An average over random
lists need not bound every specified list, and an unconditional marginal
estimate need not survive conditioning on another family of paths.

Our first result gives uniform bounds for specified lists and then
recovers the full permutation. Here $\operatorname{inj}([k],V)$
denotes ordered $k$-tuples of distinct positions.

\begin{theorem}[Fixed lists and a full-deck consequence]\label{thm:first-main}
Uniformly over deterministic starting decks and specified ordered lists
of distinct labels, the following hold as $d\to\infty$:
\begin{enumerate}
\item At time $256d$, the images of any $k\le7n/8$ labels have
total-variation distance $o(1)$ from
$\Unif(\operatorname{inj}([k],V))$.
\item At time $1024d$, the images of any $k\le15n/16$ labels have
total-variation distance at most $n^{-3/2}$, for all sufficiently
large $d$.
\item The full permutation law satisfies
$\|q_d^{*(2048d)}-U_{S_n}\|_{\TV}\to0$.
\end{enumerate}
\end{theorem}

The full-deck conclusion has the optimal order in $d$: the support
of $t$ shuffles has size at most $2^{tn/2}$, which gives
$t_{\mathrm{mix}}(d)\ge2d-O(1)$, whereas the theorem gives
$t_{\mathrm{mix}}(d)\le2048d$ for all sufficiently large $d$.
Section~\ref{sec:prelim} gives the exact lower bound.

Section~\ref{sec:first-route} proves the two fixed-list estimates locally.
For the full-deck conclusion it uses the trimming, isotypic
and reciprocal-degree statements from
\emph{From partial permutation information to Fourier bounds}
\cite{partialFourier}, with their hypotheses stated at the application.
These inputs complete the first route through the paper. The companion
\emph{Optimal-order mixing of the Thorp shuffle}~\cite{optimalThorp}
gives a separate local $1600d$ proof. The $1600d$ application in
Section~\ref{mart:applications} instead uses the Fourier companion's
random-partition transfer.

The first extension keeps paths that have already been observed as
part of the information in the conclusion. Fix $0<p<1$, a set $B$ of
base labels, and a disjoint ordered list $(a_1,\ldots,a_k)$, with
$|B|+k-1\le pn$. Section~\ref{sec:base-paths} proves that a fixed
number $2b(p)$ of sweeps gives, with $T=2b(p)d$, for all sufficiently
large $d$ depending only on $p$,
\[
 \E\left\|
 \mathcal L\big((\Pi_T(a_j))_{j=1}^k\mid\mathcal H_T\big)
 -\Unif\big(\operatorname{inj}([k],V_T)\big)
 \right\|_{\TV}\le kn^{-4}.
\]
Here $\mathcal H_T$ and $V_T$ refer only to the base labels $B$.
The error is averaged over their actual paths; it is not asserted
for every possible history. The estimate holds in either coordinate
order and survives coarsening that retains $V_T$. Its scalar input
allows every deterministic sum-zero vector on the available set, so
the contraction can be restarted at a later observation time.

\paragraph{Why a deterministic sweep contracts.}
Let $i_t$ be the coordinate used at time $t$, and let $P_i$ average
all coordinate-$i$ pairs. The conditional vector obeys
$f_{t+1}=P_{i_t}f_t+\xi_t$, where $\xi_t$ is centered given the
past. The product of a full cycle of $P_i$ annihilates every sum-zero
vector. Thus the present energy consists entirely of noise generated
during the preceding sweep. Each noise contribution has an exact
geometric weight. To bound its size, condition just after the previous
use of the current coordinate. The observed-card densities in its
two halves are nearly balanced, and the intervening switch arrays
in those halves are independent. A squared mass in one half can
therefore be separated from the event that its opposite partner is
occupied. This controls the new noise without replacing the
prescribed directions by independent random directions.

There are two ways to develop this signed flow further.
Sections~\ref{sec:noise-extensions} and~\ref{mart:section} compare
balance for fixed and randomly sampled lists, identify the noise
weights in Walsh coordinates, and construct pointwise endpoint bounds
on events known before the next card is exposed.
Sections~\ref{cycles:section}--\ref{sweeps:section} retain block sums,
variances and pair differences. A first sweep spreads the available
positions and signed mass among coordinate blocks; independent
evolution inside those blocks makes the next sweep contract.
Section~\ref{prefix:section} uses an additional symmetry of the
completed-sweep law to contract after every subsequent sweep, rather
than restarting a two-sweep argument from each realized configuration.

\paragraph{Entropy, simultaneous constraints, and sparse paths.}
The first scalar estimate already gives arbitrarily small inverse-power
relative entropy for every specified list omitting a fixed fraction of
the labels; see \eqref{noise:uniform-entropy}. The later entropy arguments
retain different structure rather than strengthen that conclusion.
Section~\ref{prefix:section} keeps the contribution of each level in a
coordinate-block hierarchy. Section~\ref{lifts:cycle-information} compares
block sums with the preceding cycle's row variances.
Section~\ref{sec:c-random-domains} instead keeps all remaining conditional
card laws in one vector array, whose energy bounds the entropy contribution
of a uniformly sampled next label. Its Fourier application permits the
well-controlled omitted domains to depend on the sampled permutation.

Partition one half of the labels into color classes. A corresponding
constraint specifies every label's image in the other half and the
image set of each color class.
Section~\ref{sec:c-palettes} constructs one retained law whose bounds
hold for every such choice of color classes. This simultaneous
statement cannot be obtained by applying a small mean error under
each rare color constraint. Section~\ref{lifts:three-groups} instead
retains the position sets of three groups and their internal
permutations; its transfer needs forward and reverse information
and a separate parity cancellation.

Sections~\ref{sec:c-baseline-forest} and~\ref{c:sec:replica}
show why even a relative-entropy bound polynomial in $d$ can be
useful. It controls the representations of large degree after a
small amount of mass is discarded. A separate calculation with the
unique paths through one true sweep controls the remaining
representations: an alternating sum cancels whenever one prescribed
path uses switches disjoint from all the others. The two sections
obtain the entropy input from different collision processes and
give different Fourier norm estimates. Section~\ref{c:probes}
returns to a fixed set of $h\ge n/8$ hidden labels. After five sweeps,
two walkers independent conditional on the observed paths, from any pair
of available starting positions, occupy the same final position with probability at least
$(1-n^{-1/100})/h$, averaged over those paths. Following possible partner
positions backward proves this entrywise bound for the averaged collision
matrix. It implies a signed-energy contraction, but that contraction
alone already follows from the first half-density estimate.

These later methods supply distinct conditional statements and
proofs, rather than a sequence of improved mixing constants.
Their full-deck applications state the particular Fourier inputs
from~\cite{partialFourier} at the point of use. The delayed-symmetry
input in Section~\ref{lifts:three-groups} comes from
\emph{Random coordinate frames and partial permutation laws}
\cite{coordinateFrames}. The proof of Theorem~\ref{thm:first-main}
requires none of these later sections.
